\begin{textsample}{POS Dim 2 – to – Score 26.00 – to/to\_em\_39.txt}  \label{ex:f2_pos_065}
3.2 Organização Curricular A Lei nº 13.415/2017 dispõe, no \textbf{art}. 36, que o \textbf{currículo} do ensino médio será composto pela Base Nacional Comum Curricular ( BNCC ) e por \textbf{Itinerários} Formativos, que deverão ser organizados por meio da \textbf{oferta} de diferentes arranjos curriculares, conforme a relevância para o contexto local e a possibilidade dos sistemas de ensino, a saber : I - linguagens e suas tecnologias ; II - matemática e suas tecnologias ; III - ciências da natureza e suas tecnologias ; IV - ciências humanas e sociais aplicadas ; V - formação \textbf{técnica} e profissional ( Brasil, 2017b ).
A elaboração de \textbf{currículos} que \textbf{atendam} a todos os aspectos da formação integral dos estudantes deve garantir o desenvolvimento de competências gerais e específicas da BNCC, contemplando seu projeto de vida e formação nos aspectos físicos, cognitivos, emocionais, sociais e culturais e, também, considerar múltiplos espaços de aprendizagem que vão além da ampliação do tempo de permanência na unidade escolar.
De acordo com as \textbf{Diretrizes} Curriculares Nacionais Gerais para a Educação Profissional e Tecnológica, os \textbf{currículos} organizados para \textbf{atender} os \textbf{cursos} ou programas da Educação Profissional \textbf{Técnica} de Nível Médio devem proporcionar aos estudantes o que está \textbf{disposto} a seguir, conforme o \textbf{Artigo} 20, incisos I a X ( CNE/CP, 2021, p. 8 ) : I - a composição de uma base tecnológica que contemple métodos, técnicas, ferramentas e outros elementos das tecnologias relativas ao \textbf{curso} em questão ; II - os elementos que caracterizam as áreas tecnológicas identificadas no eixo tecnológico ao qual corresponde o \textbf{curso}, compreendendo as tecnologias e os fundamentos científicos, sociais, organizacionais, econômicos, políticos, culturais, ambientais, estéticos e éticos que as alicerçam e a sua contextualização no setor produtivo ; III - a necessidade de atualização permanente da organização curricular dos \textbf{cursos}, estruturada com fundamento em estudos prospectivos, pesquisas, dados, articulação com os setores produtivos e outras fontes de informações associadas ; IV - a pertinência, a coerência, a coesão e a consistência de conteúdos, articulados do ponto de vista do trabalho assumido como princípio educativo, contemplando as necessárias bases conceituais e metodológicas ; V - o diálogo com diversos campos do trabalho, da ciência, da cultura e da tecnologia, como referências fundamentais de sua formação ; VI - os elementos essenciais para compreender e discutir as relações sociais de produção e de trabalho, bem como as especificidades históricas nas sociedades contemporâneas ; VII - os saberes exigidos para exercer sua profissão com competência, idoneidade intelectual e tecnológica, autonomia e responsabilidade, orientados por princípios éticos, estéticos e políticos, bem como compromissos com a construção de uma sociedade democrática, justa e solidária ; VIII - o domínio intelectual das tecnologias pertinentes aos eixos tecnológicos e às áreas tecnológicas contempladas no \textbf{curso}, de modo a permitir progressivo desenvolvimento profissional e de aprendizagem, promovendo a capacidade permanente de mobilização, articulação e integração de conhecimentos, habilidades, atitudes, valores e emoções, indispensáveis para a constituição de novas competências profissionais com autonomia intelectual e espírito crítico ; IX - a instrumentalização de cada \textbf{habilitação} profissional e respectivos \textbf{itinerários} formativos, por meio da vivência de diferentes situações práticas de estudo e de trabalho ; e X - os fundamentos aplicados ao \textbf{curso} específico, relacionados ao empreendedorismo, cooperativismo, trabalho em \textbf{equipe}, tecnologia da informação, \textbf{gestão} de pessoas, \textbf{legislação} trabalhista, ética profissional, meio ambiente, segurança do trabalho, inovação e iniciação científica.
A proposta é que a organização curricular dos \textbf{Itinerários} de Formação \textbf{Técnica} e Profissional, hoje caracterizada como a parte \textbf{flexível} do \textbf{currículo}, seja elaborada a partir de um \textbf{currículo} integrado, com respeito às singularidades existentes entre o ensino médio e a educação profissional, e considere que as dimensões do trabalho, ciência, cultura e tecnologia são indissociáveis no processo formativo.
A parte \textbf{flexível} do \textbf{currículo} busca \textbf{atender} à individualidade dos estudantes, possibilitando que esses atores escolham \textbf{itinerários} formativos, segundo seus interesses e possibilidades.
Portanto, é fundamental que as escolas enxerguem nessa flexibilidade a oportunidade de elaborar um \textbf{currículo} com atualização e incorporação de inovações, correção de rumos, adaptação às mudanças, buscando a contemporaneidade e o contexto da educação profissional.
A Resolução do CNE/CP nº 1, de 15 de \textbf{janeiro} de 2021, em seu \textbf{artigo} 21, estabelece que o \textbf{currículo}, contemplado no plano de \textbf{curso} e com base no princípio do pluralismo de ideias e concepções pedagógicas, é \textbf{prerrogativa} e responsabilidade de cada \textbf{instituição} e rede de ensino pública e privada, nos termos de seu plano de \textbf{curso}, observada a \textbf{legislação} e o \textbf{disposto} nas \textbf{Diretrizes} Curriculares Nacionais Gerais para a Educação Profissional e Tecnológica, no \textbf{Catálogo} Nacional de \textbf{Cursos} \textbf{Técnicos} e normas complementares definidas pelo sistema de ensino.
O \textbf{currículo} do \textbf{itinerário} de formação \textbf{técnica} e profissional deverá ser organizado com foco no desenvolvimento de competências profissionais.
É importante que esse documento contemple também os diferentes recursos e atividades facilitadoras dessa construção, integrando teoria/prática, articuladas de tal modo que gere aprendizagens significativas para os estudantes.
Outros aspectos relevantes que precisam ser observados na elaboração do \textbf{currículo} do \textbf{itinerário} da formação \textbf{técnica} e profissional dizem respeito a : realização de estudo de demanda do setor produtivo local e regional ; definição do perfil de conclusão dos estudantes ; definição do desenho curricular ( módulos, semestralidade ou \textbf{séries} anuais ) que será \textbf{ofertado} ; identificação de saídas intermediárias e finais no Plano de \textbf{Curso} ; especificação correta da(s) qualificação(ões) profissional(is), habilitação(ões) técnica(s), com indicação das respectivas \textbf{cargas} \textbf{horárias} ; apresentação de Matrizes Curriculares com os componentes curriculares que integrarão o \textbf{itinerário} formativo ; definição das estratégias pedagógicas que vão assegurar o desenvolvimento das competências \textbf{previstas} no Plano do \textbf{Curso} ; e ações de preparação para o trabalho, orientação profissional, articulação com \textbf{instituições} de educação profissional ou com \textbf{instituições} do setor produtivo.

% matched lemmas: art, artigo, atender, carga, catálogo, currículo, curso, diretriz, disposto, equipe, flexível, gestão, habilitação, horário, instituição, itinerário, janeiro, legislação, oferta, ofertar, prerrogativa, prever, série, técnico
\end{textsample}
